\documentclass[aspectratio=169]{beamer}
\input{header}
\coursecode{JKSSB}

\title{Organization of a Computer}
\subtitle{Lesson 04 --- Units, Buses and the Instruction Cycle}
\author{JKSSB Computer Awareness}
\date{\today}

\begin{document}
\frame{\titlepage}

\begin{frame}{Video Lesson 04}
  \centering
  \Large\href{https://youtu.be/92XQ_n1t_cc}{Watch: How Computers Actually Work: CPU, Memory \& I/O Explained}
\end{frame}

\begin{frame}{The Office Desktop as a System}
  \begin{itemize}
    \item The clerk types marks on the \key{keyboard} and clicks with the
      \key{mouse}; the scanner feeds a form.
    \item The \key{CPU} processes the marks using a spreadsheet formula.
    \item The \key{monitor} shows the total and the \key{printer} prints it.
    \item \key{RAM} (Random Access Memory) holds the work while it runs; the
      \key{SSD} keeps the file after the machine is switched off.
  \end{itemize}
  \vspace{0.3cm}
  \muted{One familiar machine carries every unit in this lesson.}
\end{frame}

\begin{frame}{Five Functional Units}
  \begin{itemize}
    \item \key{Input unit}: accepts data and instructions.
    \item \key{Output unit}: presents the result.
    \item \key{Memory unit}: primary storage that holds data and instructions
      currently in use.
    \item \key{Arithmetic Logic Unit (ALU)}: performs arithmetic and logical
      operations.
    \item \key{Control unit}: directs the sequence of operations and the flow
      of data.
  \end{itemize}
  \begin{block}{High-yield grouping}
    The ALU and the control unit, together with registers, form the
    \key{Central Processing Unit (CPU)}.
  \end{block}
\end{frame}

\begin{frame}{Inside the CPU}
  \begin{columns}[T]
    \column{0.46\textwidth}
      \textbf{Arithmetic Logic Unit}
      \begin{itemize}
        \item Adds, subtracts, compares
        \item Works on binary data
        \item Produces a result and status flags
      \end{itemize}
    \column{0.46\textwidth}
      \textbf{Control Unit}
      \begin{itemize}
        \item Fetches and decodes instructions
        \item Sends control signals
        \item Coordinates all other units
      \end{itemize}
  \end{columns}
  \vspace{0.3cm}
  \muted{Registers are small, fast storage inside the CPU; the program counter
  and instruction register are examples.}
\end{frame}

\begin{frame}{Input and Output Units}
  \begin{columns}[T]
    \column{0.46\textwidth}
      \textbf{Input devices}
      \begin{itemize}
        \item Keyboard, mouse
        \item Scanner, microphone
        \item Light pen, digitizer
      \end{itemize}
    \column{0.46\textwidth}
      \textbf{Output devices}
      \begin{itemize}
        \item Monitor, printer
        \item Speaker, plotter
        \item Sound card output
      \end{itemize}
  \end{columns}
  \vspace{0.3cm}
  \muted{Classify the device by the role it plays in the described task.}
\end{frame}

\begin{frame}{Memory (Primary) versus Storage (Secondary)}
  \begin{block}{Primary storage / main memory}
    RAM and ROM --- directly accessed by the CPU; holds data and instructions
    currently being processed. RAM is volatile.
  \end{block}
  \begin{alertblock}{Secondary storage}
    Hard disk, SSD and pen drive --- keeps files permanently, but is not
    accessed directly during processing.
  \end{alertblock}
  \vspace{0.2cm}
  \muted{The CPU works with main memory; secondary storage is reached through
  main memory and the operating system.}
\end{frame}

\begin{frame}{The Three Buses}
  \begin{center}
    \begin{tabular}{p{0.22\textwidth} p{0.24\textwidth} p{0.44\textwidth}}
      \toprule
      \textbf{Bus} & \textbf{Direction} & \textbf{Carries} \\
      \midrule
      Address bus & CPU $\rightarrow$ memory & The address of a memory
      location \\
      Data bus & Bidirectional & Data and instructions between CPU and
      memory \\
      Control bus & CPU $\rightarrow$ units & Control and timing signals \\
      \bottomrule
    \end{tabular}
  \end{center}
  \vspace{0.3cm}
  \muted{Address = where; data = what; control = when and how.}
\end{frame}

\begin{frame}{Address vs Data vs Control Bus}
  \begin{itemize}
    \item \key{Address bus}: one-way, from the CPU to memory --- it carries
      the location being selected, not the contents.
    \item \key{Data bus}: two-way --- it moves the actual data and
      instructions in both directions.
    \item \key{Control bus}: one-way control path --- it carries read, write
      and timing signals.
    \end{itemize}
  \begin{alertblock}{Near-neighbour trap}
    The address bus does \emph{not} carry data, and the data bus does
    \emph{not} carry addresses. Do not swap them.
  \end{alertblock}
\end{frame}

\begin{frame}{Motherboard, Slots and Interfaces}
  \begin{itemize}
    \item The \key{motherboard} is the main circuit board linking the CPU,
      memory and buses.
    \item \key{Expansion slots} let you add cards such as a graphics card or
      a network card.
    \item \key{Interfaces/ports} connect external devices such as the
      keyboard, printer and monitor.
    \item The system bus runs along the motherboard between the CPU and the
      other units.
  \end{itemize}
\end{frame}

\begin{frame}{The Stored-Program (von Neumann) Idea}
  \begin{itemize}
    \item Early machines were rewired for each task.
    \item The \key{stored-program concept} keeps the program's instructions
      in the same main memory as the data.
    \item The CPU fetches an instruction, then processes data under it ---
      no rewiring needed.
  \end{itemize}
  \begin{block}{Association}
    The idea that both instructions and data are stored in memory is
    associated with \key{von Neumann architecture}.
  \end{block}
\end{frame}

\begin{frame}{The Instruction Cycle}
  \begin{enumerate}
    \item \key{Fetch}: the CPU obtains the next instruction from main
      memory.
    \item \key{Decode}: the control unit works out what the instruction
      means.
    \item \key{Execute}: the ALU or another unit performs the operation.
    \item \key{Store} (when required): the result is written back to memory
      or a register.
  \end{enumerate}
  \begin{alertblock}{Order cue}
    The cycle repeats in the fixed order fetch $\rightarrow$ decode
    $\rightarrow$ execute.
  \end{alertblock}
\end{frame}

\begin{frame}{Data and Instruction Flow}
  \begin{itemize}
    \item Input data and instructions enter through input devices.
    \item They are held in main memory and fetched by the CPU over the
      system bus.
    \item The CPU decodes and executes, using the ALU and registers.
    \item Results go to output devices or back to main memory.
    \item Secondary storage keeps files between sessions --- the CPU does
      not reach it directly.
  \end{itemize}
\end{frame}

\begin{frame}{Lesson 04: Recall}
  \begin{itemize}
    \item Five units: input, output, memory, ALU, control unit;
      CPU = ALU + control unit + registers.
    \item Address bus: CPU $\rightarrow$ memory, carries the address. Data
      bus: bidirectional, carries data and instructions. Control bus: carries
      control signals.
    \item The CPU directly accesses main memory (primary storage), not
      secondary storage.
    \item Stored-program (von Neumann): instructions and data share memory.
    \item Instruction cycle: fetch $\rightarrow$ decode $\rightarrow$ execute.
  \end{itemize}
\end{frame}
\end{document}
